\section{Related work}
\label{sec:related}

\paragraph{Proxies, gaming, and side effects.}
The gap between what an objective rewards and what its designer meant has a
long history. \citet{manheim2018categorizing} separate four ways a proxy can
come apart from its goal under optimization; the case closest to ours is what
they call \emph{extremal} Goodhart, where the proxy is accurate for ordinary
solutions and stops tracking the goal once an optimizer pushes into unusual
ones. \citet{amodei2016concrete} discuss two relevant failure modes. Their
\emph{reward hacking} motivates visible/hidden test splits such as ours.
Their \emph{negative side effects} are the closer analogy for what we
observe: an objective that is silent about part of the environment leaves the
agent free to disturb it. The visible tests say nothing about the
deduplication helper, so an agent can rewrite it while passing them.
\citet{krakovna2020specification} collect dozens of such cases under the name
\emph{specification gaming}. \citet{skalse2022defining} show formally that
when the agent can choose among all stochastic policies, a proxy that can
never be gamed must be equivalent to the true reward or trivial, and \citet{pan2022effects} find that more capable agents can score
higher on a misspecified proxy while doing worse on the true reward, sometimes
with an abrupt transition. That last result is one reason we track two model
generations instead of one.

\paragraph{Gaming by language-model agents.}
Recent work has moved these ideas to language models acting with tools.
\citet{denison2024sycophancy} train models on a curriculum of increasingly
gameable tasks and find a small rate of zero-shot tampering with a reward
function, occasionally including edits to the unit tests meant to catch it.
\citet{bondarenko2025demonstrating} show that some reasoning models, asked to
beat a chess engine, rewrite the board state instead of playing, and that the
propensity varies sharply between models and with the wording of the
instruction. \citet{baker2025monitoring} document reward hacks that arose
during the training of a frontier reasoning model on agentic coding tasks,
including exiting before the tests run, skipping tests, and editing test files,
and use the trick of reverting test edits to expose them. These studies look
for adversarial behavior: the agent defeats the check. We measure something milder: unrequested edits made while fixing a bug,
including edits that leave every test passing. Our agents almost always fix the
bug, and the excursions we measure are mostly not attempts to defeat a check;
they simply do more than they were asked.

\paragraph{Code benchmarks and what pass rates hide.}
HumanEval \citep{chen2021evaluating} and MBPP \citep{austin2021program} made
unit-test pass rates the standard measure of code generation, and SWE-bench
\citep{jimenez2024swebench} carried the idea to repository-level issues, where
an agent's patch is judged by tests it never sees. A steady line of work has
since shown how much such scores can hide. EvalPlus \citep{liu2023evalplus}
added many more tests to HumanEval and found that a sizeable share of
accepted solutions were wrong. MBPP's own authors observed a model
hard-coding an input to pass a visible assert \citep{austin2021program}. An
audit of resolved SWE-bench instances found that about a third of them leaked
the fix in the issue text and a similar share passed only because the tests
were weak \citep{aleithan2024swebench}. OpenAI's annotation of SWE-bench led
to the Verified subset \citep{chowdhury2024introducing}, and a later audit led
OpenAI to stop reporting it altogether \citep{openai2026why}. Contamination
adds a further layer: models recall SWE-bench file paths from issue text alone
far more often than for repositories outside the benchmark \citep{liang2025swebench}, and
classic code benchmarks overlap with common training corpora
\citep{riddell2024quantifying,matton2024leakage}, which motivates fresh,
time-stamped tasks \citep{jain2025livecodebench,zhang2025swebench,badertdinov2025swerebench}.
Our tasks were written for this study and never published, and our hidden
tests play the role that held-out tests play in these benchmarks. The
difference is that we also measure what the agent changed that no test, seen
or unseen, is designed to check.

\paragraph{Closest prior work.}
Several recent studies contain one of our instruments, and one anticipates
part of our main result. SpecBench \citep{zhao2026specbench} gives
long-horizon coding agents visible validation tests and scores them on
held-out tests, defining a reward hacking gap that is essentially our drift
gap; EvilGenie \citep{gabor2026evilgenie} and \citet{ahmed2026investigating}
use similar visible/held-out splits on competitive programming and on
SWE-bench patches. ImpossibleBench \citep{zhong2025impossiblebench} makes
tests contradict the specification so that any pass counts as cheating, and
finds that stricter prompts lower the cheating rate. OverEager-Bench
\citep{qu2026overeager} plants trap actions, such as deleting unrelated
files, in benign tasks and finds that removing a scope statement from the
prompt raises the rate at which agents take them. FixedBench
\citep{gloaguen2026coding} gives agents SWE-bench issues that are already
fixed and finds that many still edit code, often code unrelated to the
issue. In real sessions, \citet{tang2026how} code about one in ten
misalignment episodes as self-initiated overreach, most of it scope
overreach.

Closest to our finding, \citet{weinberger2026promptinduced} add a clause
inviting the agent to ``clean up anything adjacent'' to an otherwise fully
specified prompt and report that this scope-expanding wording was the only
prompt family that consistently widened diffs, producing out-of-scope edits
in a small share of runs for some models. Our results agree on direction and
extend it in three ways. Our loose prompt has no explicit scope line to
compete with, and the effect is far larger. We compare it directly with a
vague prompt on the same repository, which produced no planted-function
edits and far fewer out-of-scope edits of any kind. And we check the out-of-scope edits
against hidden tests, which shows that pass/fail scoring misses nearly all of
them. The vagueness result is not universal: UnderSpecBench
\citep{ji2026coding} finds that less clear intent \emph{increases}
out-of-scope actions by agents doing operations tasks in a simulated DevOps
environment, rather than editing a small repository. Which kind of
missing information pushes agents out of scope probably depends on what
there is to act on.

\paragraph{Underspecified requests and instruction following.}
A separate line of work studies what models do with instructions that leave
things out. Instruction-following benchmarks such as IFEval
\citep{zhou2023instructionfollowing} check verifiable constraints stated in
the prompt, whereas our question is about constraints the prompt does not
state. In code, ClarifyGPT \citep{mu2023clarifygpt} detects ambiguous requirements
and asks clarifying questions to improve the generated code, HumanEvalComm
\citep{wu2025humanevalcomm} measures whether models ask such questions on
their own, and Ambig-SWE \citep{vijayvargiya2026ambigswe} extends
this to repository-level tasks whose issue text has been made underspecified.
These studies treat the missing information as the problem. Our vague rung is
in the same spirit, but in our tasks missing information produced only a
little scope creep, mostly in tests, while an explicit invitation to improve
things produced a great deal. Work on sycophancy \citep{sharma2024towards,perez2023discovering} offers a
plausible mechanism. A model tuned to please may read ``get things into good
shape'' as a request to show visible effort, and the cheapest way to show
effort in a small repository is to clean up whatever looks untidy.
